\documentclass[12pt,a4paper]{article}
\usepackage[utf8]{inputenc}
\usepackage{amsmath, amssymb, bm}
\usepackage{algorithm}
\usepackage{algorithmic}
\usepackage{graphicx}
\usepackage{geometry}
\usepackage{caption}
\usepackage{cite}
\usepackage{listings}
\usepackage{xcolor}
\usepackage{tikz}
\usetikzlibrary{positioning, shapes.geometric, arrows.meta, fit, backgrounds, calc}
\geometry{margin=1in}

% 启用中文支持
\usepackage{ctex}

% 放宽 hbox 行末对齐以消除长标识符造成的 overfull
\emergencystretch=2em
\tolerance=2000
\hbadness=5000

% 长 typewriter 标识符的细粒度换行辅助宏
\newcommand{\brk}{\discretionary{}{}{}}

% listings 代码块配置
\lstset{
    basicstyle=\ttfamily\small,
    keywordstyle=\color{blue!70!black}\bfseries,
    commentstyle=\color{gray!70!black}\itshape,
    stringstyle=\color{orange!80!black},
    numbers=left,
    numberstyle=\tiny\color{gray},
    frame=single,
    breaklines=true,
    showstringspaces=false,
    language=Python,
}

% 将 hyperref 放置在最后，并显式配置内部链接和颜色
\usepackage[colorlinks=true, linkcolor=blue, anchorcolor=blue, citecolor=blue, urlcolor=blue, bookmarksnumbered=true]{hyperref}

\title{\textbf{KABO：一种任务解耦的知识增强贝叶斯优化框架 \\ 及其在光催化 CO$_2$ 还原中的应用}}
\author{Jiayu Wu}
\date{\today}

\begin{document}
\maketitle

\begin{abstract}
仿酶催化体系（CO$_2$RR、OER、NRR、ORR 等）的实验寻优长期面临高维设计空间、昂贵实验成本和领域知识难以结构化注入三重障碍。本文提出并实现了一种\textbf{任务解耦的知识增强贝叶斯优化框架}（Task-decoupled Knowledge-Augmented Bayesian Optimization, KABO），将纯数学优化能力封装为与体系无关的 \texttt{KABOEngine}，将体系相关逻辑抽象为统一的 \texttt{TaskBase} 接口。算法核心整合了偏好学习（Preference Learning）、专家先验编码（Expert Prior Encoding）、偏好探索子循环（Preference Exploration）与近似信息价值（Approximate VOI）四项知识内化机制，通过 z-score 标定的复合采集函数协同驱动候选选择。为贴合真实催化体系中"连续-整数-类别"共存的描述符现实，进一步实现了\textbf{混合离散-连续搜索空间支持}：Task 层通过 \texttt{feature\_types()} 契约声明维度类型，算法层以整数 Round-trick、\texttt{MixedSingleTaskGP} 混合核与 Thompson 离散策略三位一体实现类型感知的代理建模与采集优化。任务层以一个体系一个文件的形式组织（\texttt{kabo/task/*.py}），已内置光催化 CO$_2$ 还原任务（\texttt{CO2RRTask}，19 描述符 / 7 产物，其中 6 维为整数氢键计数）与最小综合任务（\texttt{TestTask}，3 描述符 / 1 产物）。新增体系（OER 等）仅需新增一个约 50 行的 Task 类，Engine、Optimizer 与 CLI 零改动。本文以 CO$_2$RR 为主体案例，详述了架构分层、数学重构、KABO 四项机制、混合空间处理、偏好探索子循环与人类在环闭环；并通过 TestTask 的端到端冒烟测试印证 Engine 的体系无关性。

\textbf{关键词}：贝叶斯优化；高斯过程；知识增强；偏好学习；任务解耦架构；光催化；CO$_2$ 还原
\end{abstract}

\vspace{2em}
\tableofcontents
\newpage


\section{引言 (Introduction)}

仿酶催化体系——以光催化 CO$_2$ 还原 (CO$_2$RR) 为代表——在通往化学品绿色合成的路径上具有战略价值\cite{nature2025, jacs2024}，但其材料探索在工程上面临三重结构性困难：

\begin{enumerate}
    \item \textbf{高维非凸设计空间}：催化剂性能受数十个描述符共同控制（能带间隙、A\_pI、位点距离、溶剂偶极……），配方响应面高度非线性，全穷举与高通量筛选 (HTS)\cite{nature2025} 受制于"维数灾难"。
    \item \textbf{昂贵实验评估函数}：单次光催化收率测量通常需要数小时到数天，实验预算难以支撑样本数 $N \gg$ 维度 $d$ 的经典统计假设。
    \item \textbf{领域知识的结构化注入缺失}：资深化学家的经验——"最优电位约 $-0.8$ V"、"甲酸产物的选择性与亲水溶剂强相关"——长期停留在口头约定与手工试错，难以融入数学模型结构。
\end{enumerate}

应对昂贵评估的经典范式是贝叶斯优化 (Bayesian Optimization, BO)\cite{bo2012}，其在黑盒函数导数不可得的条件下提供了理论严密的主动采样策略\cite{botorch2019, turbo2019}。然而标准 BO 与真实实验室之间仍存在两类鸿沟：其一，量纲悬殊导致核函数的欧氏度量坍塌\cite{turbo2019}；其二，专家直觉无法被代理模型结构性地"内化"。本报告在原始"人类在环"贝叶斯优化 (HITL-BO) 框架基础上，构建了 \textbf{KABO 知识增强贝叶斯优化}体系，并通过一层任务抽象使其天然支持多类催化体系。

KABO 将优化流程组织为四阶段：降维筛选 (Phase 1) $\rightarrow$ 高斯代理模型重建 (Phase 2) $\rightarrow$ 偏好探索与知识采集 (Phase 2.5) $\rightarrow$ 知识增强采集-交互闭环 (Phase 3)，并通过偏好学习\cite{astudillo2019}、专家先验编码\cite{li2020prior}及近似信息价值\cite{wu2016kg}三条路径实现知识的持久内化与循环利用。Task 抽象层则负责所有体系相关语义（特征、产物、边界、交互文案），使同一份算法内核可以驱动任意催化体系。


\section{任务解耦的系统架构 (Task-Decoupled System Architecture)}

\subsection{分层职责}

KABO 采用三层架构（图 \ref{fig:arch}）将纯数学能力与体系语义彻底解耦：

\begin{itemize}
    \item \textbf{算法层} \texttt{kabo.engine.KABOEngine}：聚合 \texttt{SurrogateModel} (GP 代理)、\texttt{PreferenceModel} (PairwiseGP 偏好学习)、\texttt{ExpertPrior} (专家先验打分)。接口仅接受数值矩阵、边界张量与特征列表，\textbf{不包含任何体系关键词}。
    \item \textbf{领域层} \texttt{kabo.task.TaskBase}：抽象合同定义特征集合、设计域边界、产物列、训练目标构造、交互与模拟逻辑。每个体系以一个独立 Python 文件实现（\texttt{kabo/task/co2rr.py}、\texttt{kabo/task/test\_task.py} 等），通过 \texttt{@register\_task} 装饰器自动注册到全局任务表 \texttt{TASK\_REGISTRY}。
    \item \textbf{编排层} \texttt{kabo.optimizer.KABOOptimizer}：负责三阶段流程调度、离散/连续候选合并、Top-N 菜单呈现与审计字段写入。编排器持有一个 \texttt{task} 对象，所有体系相关读操作均经由 \texttt{self.task.*} 委派。
\end{itemize}

\begin{figure}[htbp]
\centering
\footnotesize
\begin{tikzpicture}[
    font=\footnotesize,
    node distance=0.9cm and 0.8cm,
    every node/.style={align=center},
    layerlabel/.style={font=\footnotesize\bfseries\itshape, text=gray!60!black},
    box/.style={
        draw=black!70, thick, rounded corners=3pt,
        fill=white, inner sep=6pt, minimum height=1cm,
    },
    cli/.style   ={box, fill=blue!8,   text width=9cm},
    orch/.style  ={box, fill=green!8,  text width=9cm},
    engine/.style={box, fill=orange!8, text width=6.4cm},
    task/.style  ={box, fill=purple!8, text width=6.4cm},
    arrow/.style={-{Stealth[length=2.5mm]}, thick, black!70},
  ]

  \node[cli] (cli) {%
    \textbf{CLI / run.py}\quad\texttt{(kabo/cli.py)}\\[1pt]
    解析 \texttt{-{}-task} / \texttt{-{}-data} / \texttt{-{}-iterations}\\
    从 \texttt{TASK\_REGISTRY} 构造 Task
  };

  \node[orch, below=of cli] (orch) {%
    \textbf{KABOOptimizer}\quad\texttt{(kabo/optimizer.py)}\\[1pt]
    Phase 1 特征选择 (RF)\ \textbar{}\ Phase 2 GP 拟合\ \textbar{}\ Phase 3 采集 + 人机交互循环
  };

  \node[engine, below=1.2cm of orch.south, xshift=-3.4cm, anchor=north] (eng) {%
    \textbf{KABOEngine}\quad\texttt{(kabo/engine.py)}\\[1pt]
    SurrogateModel \textbullet\ PreferenceModel \textbullet\ ExpertPrior\\
    UCB / qNEI / KABO\\[1pt]
    \textcolor{red!70!black}{\small$\triangle$ 无体系关键词}
  };

  \node[task, below=1.2cm of orch.south, xshift=3.4cm, anchor=north] (task) {%
    \textbf{TaskBase (子类)}\quad\texttt{(kabo/task/*.py)}\\[1pt]
    特征 \& 边界 \textbullet\ 产物列 \& 目标\\
    交互 \& 模拟 \textbullet\ 先验 schema
  };

  \draw[arrow] (cli.south) -- (orch.north);
  \draw[arrow] (orch.south) -- ($(eng.north)+(0,0)$);
  \draw[arrow] (orch.south) -- ($(task.north)+(0,0)$);

  \node[layerlabel, right=0.15cm of cli.east]  {装配层};
  \node[layerlabel, right=0.15cm of orch.east] {编排层};
  \node[layerlabel] at ($(eng.south)+(0,-0.3)$)  {算法层};
  \node[layerlabel] at ($(task.south)+(0,-0.3)$) {领域层};
\end{tikzpicture}
\caption{KABO 三层架构：装配、编排、算法、领域各司其职。}
\label{fig:arch}
\end{figure}

\subsection{TaskBase 接口契约}

\texttt{TaskBase} 作为抽象基类定义了以下最小契约：

\begin{itemize}
    \item \texttt{task\_name()} $\rightarrow$ 体系短标识（\texttt{"CO2RR"}、\texttt{"OER"}）
    \item \texttt{feature\_columns()} $\rightarrow$ 全量特征列名（有序）
    \item \texttt{design\_space\_bounds()} $\rightarrow$ 每个特征的物理边界
    \item \texttt{target\_columns()} $\rightarrow$ 产物短名 $\rightarrow$ DataFrame 列名映射
    \item \texttt{default\_target()} $\rightarrow$ 默认优化目标
    \item \texttt{build\_training\_target(df, target\_column, **kwargs)} $\rightarrow$ 代理模型训练向量构造（支持复合目标，如 CO$_2$RR 的 $Y_{\text{target}} - w \cdot Y_{H_2}$）
    \item \texttt{prompt\_observation(target\_column)} $\rightarrow$ 交互模式下收集实验观测
    \item \texttt{simulate\_observation(target\_column, $\mu$, $\sigma$)} $\rightarrow$ 演示模式下模拟观测
\end{itemize}

可选 Hook（\texttt{validate\_candidates}、\texttt{parse\_expert\_prior}）提供默认实现。契约设计遵循"已知抽象、未知具体"原则：Engine 只向 Task 提出问题（"你的特征是什么？目标是什么？"），Task 的具体答案不渗漏回算法层。

\subsection{注册机制与零改动扩展}

新增体系的工作流具体化为三步：

\begin{lstlisting}[caption={新增 OERTask 的最小模板},label={lst:oer}]
# kabo/task/oer.py
from kabo.task.base import TaskBase, register_task

@register_task
class OERTask(TaskBase):
    def task_name(self): return "OER"
    def feature_columns(self): return ["overpotential", "pH", ...]
    def design_space_bounds(self):
        return {"overpotential": (0.2, 1.2), "pH": (0.0, 14.0), ...}
    def target_columns(self): return {"O2": "Y_O2"}
    def default_target(self): return "O2"
    def build_training_target(self, df, target_column, **kwargs):
        return df[target_column].values.astype(float)
    def prompt_observation(self, target_column): ...
    def simulate_observation(self, target_column, y_mean, y_std): ...
\end{lstlisting}

随后在 \texttt{kabo/task/\_\_init\_\_.py} 追加一行 \texttt{from kabo.task.oer import OERTask} 以触发注册。用户立即可用 \texttt{python run.py -{}-task oer}，\textbf{Engine、Optimizer、CLI 核心代码零改动}。这一约束由算法层模块的静态代码检查保证：对 \texttt{kabo/engine.\brk py}、\texttt{kabo/acquisition.\brk py}、\texttt{kabo/feature\_\brk selection.py} 等算法层文件执行 \texttt{grep} 扫描，不应出现任何体系特定常量。


\section{描述符处理与边界感知归一化}

\subsection{非线性特征拆分及基尼重要度评判}

原始数据集在 Task 声明的 $d$ 维描述符空间中展开（CO$_2$RR 为 19 维）。在可获取数据点数量有限（$N \sim 20$）、局部信噪比低的条件下强行使用全维 GP 易导致过度平滑或局部震荡。系统前置随机森林非线性降维，使用基尼重要度保留 Top-$K$ 维：

\begin{equation}
    \text{Imp}(x_j) = \frac{1}{N_{trees}} \sum_{t=1}^{N_{trees}} \sum_{s \in \text{nodes}(t), v(s) = x_j} p(s) \cdot \Delta I(s)
\end{equation}

$\Delta I(s)$ 为基于特征 $x_j$ 拆解引起的杂质（方差）下降，$p(s)$ 为节点权重。随机森林在 $N < 10$ 时自动降低 \texttt{n\_estimators} 与深度以避免过拟合，生成 $X_{\text{selected}} \in \mathbb{R}^{N \times K}$。相较于 PCA 带来的物理解释性丧失，基尼度量保留了特征的可追溯性\cite{turbo2019}。

\subsection{缺失特征重构与"Pre-fill-before-choice"机制}

得到最优连续点 $x^*_{\text{cont}} \in \mathbb{R}^K$ 后，实验配方需恢复到 Task 声明的全部 $d$ 维以满足合成可行性。架构引入 \textbf{Pre-fill-before-choice} 机制：未被选中的 $d-K$ 维特征按专家预填值或设计域中点补齐，使候选呈现给专家时已为完整配方，避免描述符不完整导致的审核错位。

\subsection{边界感知归一化 (Bounds-Aware Normalization)}

传统机器学习常用 $\text{MinMax}(X_{\text{train}})$ 归一化，但当模型探索超出训练样本极值范围时会产生灾难性失真。系统改用 Task 提供的理论设计域边界 $(l^{(k)}, u^{(k)})$ 做归一化：

\begin{equation}
    x_{\text{norm}, k} = \max \left( 0, \;\min \left( 1, \; \frac{x_k - l^{(k)}}{u^{(k)} - l^{(k)}} \right) \right)
\end{equation}

此举确保高斯代理能够在训练样本未覆盖的区域无缝外推，完全杜绝由训练集极值作为归一化分母引入的伪外推偏误。


\section{高斯过程代理模型}

\subsection{ARD Mat\'{e}rn 核}

利用观测 $\mathcal{D}_N = \{X_{\text{selected}}^{(i)}, y^{(i)}\}_{i=1}^N$ 构建代理模型 $f(x) \sim \mathcal{GP}(\mu, k(x, x'))$。化学体系内诸如功函数 (eV) 与比表面积 ($\text{m}^2/\text{g}$) 的标度差异悬殊，通用欧氏度量会坍塌。选择带\textbf{自动关联性决定 (ARD)} 的 Mat\'{e}rn 核，$\nu = 2.5$ 兼顾不可导的高频震荡：

\begin{equation}
    k_{\text{Mat\'{e}rn } 2.5}(x, x') = \sigma_f^2 \left( 1 + \sqrt{5} d_{\text{ARD}} + \tfrac{5}{3} d_{\text{ARD}}^2 \right) \exp\left( -\sqrt{5} d_{\text{ARD}} \right)
\end{equation}

加权对角矩阵对 $K$ 维变量实现独立长度尺度缩放 $l_k$：

\begin{equation}
    d_{\text{ARD}}(x, x') = \sqrt{ \sum_{k=1}^K \frac{(x_k - x'_k)^2}{l_k^2} }
\end{equation}

$l_k$ 的自适应对特征量纲刚性错位实现有效补偿。对于更为复杂的形貌突变响应曲面，复合谱混合核 (Spectral Mixture Kernels)\cite{catbox2025} 作为可选接口通过 \texttt{--kernel-type spectral\_mixture} 启用。

\subsection{超参数的边际对数似然估计}

核方差 $\sigma_f^2$、观测白噪声方差 $\sigma_n^2$ 与尺度向量 $\mathbf{l}$ 通过最大化 Exact Marginal Log-Likelihood 获得：

\begin{equation}
    \log p(\mathbf{y} | X, \bm{\theta}) = -\tfrac{1}{2} \mathbf{y}^\top (K_\theta + \sigma_n^2 I)^{-1} \mathbf{y} - \tfrac{1}{2} \log|K_\theta + \sigma_n^2 I| - \tfrac{N}{2} \log 2\pi
\end{equation}

求解过程使用 L-BFGS-B 或基于 Torch 的梯度下降（\texttt{fit\_gpytorch\_mll}）。


\section{基础采集策略}

高分辨收率测定常带有光源扰动或杂质吸附引入的环境噪声。系统融合了置信区间探索与带噪期望改进两类算法基准，统一封装在 \texttt{KABOEngine.\brk build\_acquisition(beta)} 接口下，通过 \texttt{-{}-acq-strategy} 切换。

\subsection{UCB 与 $\beta_t$ 调度}

依 Ding et al.\cite{gpucb2026} 的 GP-UCB 理论准则，最大化置信上边界在探索与收割之间给出黄金折中。本实现采用 BoTorch 常见参数化：

\begin{equation}
    \alpha_{\text{UCB}}(x) = \mu_{t-1}(x) + \beta_t \, \sigma_{t-1}(x)
\end{equation}

三种 $\beta_t$ 调度模式：

\begin{itemize}
    \item \texttt{fixed}：常数 $\beta$，默认 2.0
    \item \texttt{theory}：$\beta_t = \beta \cdot 2 \log \left( t^{d/2 + 2} \pi^2 / (3 \delta) \right)$，允许用户 \texttt{beta} 缩放
    \item \texttt{theory-strict}：纯理论 $\beta_t$（不做缩放），用于高保真论文复现
\end{itemize}

\subsection{噪声下的 qNEI}

BoTorch 引擎底层的带噪期望提升（\texttt{qNoisyExpectedImprovement}）用 Monte Carlo 法整合观测概率提升\cite{cbo2017}：

\begin{equation}
    \alpha_{\text{qNEI}}(x) = \mathbb{E}_{y|x \sim \mathcal{N}}\left[ \max(y - y_{\text{best}}^*,\, 0) \right]
\end{equation}

使用 128 采样数（可调）的 \texttt{SobolQMCNormalSampler} 做 Quasi-Monte Carlo 无偏估计\cite{botorch2019}。qNEI 对高噪声场景比 UCB 更稳健。


\section{混合离散-连续搜索空间支持}
\label{sec:mixed-space}

\subsection{动机：催化描述符并非全连续}

仿酶催化的真实描述符空间并非纯连续：氨基酸 / 溶剂的氢键供受体数天然为整数计数 ($0, 1, \ldots, 6$)，溶剂身份、金属中心、氨基酸种类则是无序类别变量；若一律按连续变量处理，代理 GP 会在整数维上做连续插值、采集优化可能推荐 $A\_{hbond\_donors} = 2.34$ 这类在实验室不可实现的配方\cite{garrido2020round}。本节在保持向后兼容的前提下，通过\textbf{任务层类型声明 $\rightarrow$ 算法层分派}的最小侵入设计实现了混合空间支持，核心改动集中在 \texttt{kabo/task/base.py}、\texttt{kabo/surrogate.py}、\texttt{kabo/acquisition.py}、\texttt{kabo/engine.py} 四个文件。

\subsection{变量类型谱系}

为避免"一刀切"的次优处理，将催化描述符划分为五类，并为每一类配置合适的数学策略：

\begin{itemize}
    \item \textbf{连续}（温度、电位、浓度）$\to$ ARD Mat\'{e}rn GP（沿用 §4 方案）
    \item \textbf{整数}（氢键供受体数、配位数、碳链长度 $n$）$\to$ 归一化网格 Round-trick
    \item \textbf{类别无序}（溶剂身份、金属中心、氨基酸）$\to$ \texttt{MixedSingleTaskGP} 的 Hamming 核\cite{ru2020cocabo}
    \item \textbf{类别有序}（粒径等级、结晶度档）$\to$ Ordinal 核 / 单调 warp（预留）
    \item \textbf{离散数值}（预定浓度档位）$\to$ \texttt{fixed\_features\_list} 枚举（预留）
\end{itemize}

CO$_2$RR 的 19 描述符中 6 维为整数（氨基酸 A / B / 溶剂各 2 个氢键计数），13 维为连续；氨基酸身份目前以连续描述符（pI、距离）间接编码，留作 P4 的化学嵌入扩展方向。

\subsection{任务层类型声明契约}

\texttt{TaskBase} 新增两个可选方法形成"类型契约"，默认实现保持全连续行为从而不破坏历史 Task：

\begin{itemize}
    \item \texttt{feature\_types()} $\rightarrow$ \texttt{\{feature\_name: "continuous" | "integer" | "categorical" | "ordinal"\}} 字典
    \item \texttt{categorical\_values()} $\rightarrow$ 类别变量的合法取值列表
\end{itemize}

\texttt{CO2RRTask} 覆写 \texttt{feature\_types()}，显式标注 6 个 \texttt{*\_hbond\_acceptors / donors} 为 \texttt{integer}。此举使\textbf{"氢键计数是整数"这一体系知识归属 Task 层}，Engine 仅按类型字符串分派，保持"零体系关键词"原则不变。

\subsection{整数变量：归一化网格 Round-trick}

依 Garrido-Merch\'{a}n \& Hern\'{a}ndez-Lobato\cite{garrido2020round} 的 "round within GP" 策略，整数维在归一化 $[0, 1]^K$ 空间上\textbf{仍做连续优化}，但优化完成后 snap 到合法整数网格。设整数维 $k$ 的原始边界为 $(l_k, h_k)$，其归一化合法网格为
\begin{equation}
    \mathcal{G}_k = \{0, 1/n_k, 2/n_k, \ldots, 1\}, \quad n_k = h_k - l_k
\end{equation}
对最优候选 $x^*_k$ 执行
\begin{equation}
    x^*_k \leftarrow \frac{1}{n_k} \cdot \mathrm{round}(n_k \cdot x^*_k), \quad x^*_k \in [0, 1]
\end{equation}
然后在 snap 后的点上\textbf{重算采集值}，确保与离散候选池按同一尺度 apples-to-apples 排序。实现上封装为 \texttt{kabo.utils.round\_integer\_\brk dims\_\brk to\_\brk grid}，由 \texttt{optimize\_continuous} 通过 \texttt{integer\_indices} 参数驱动。

\subsection{类别变量：MixedSingleTaskGP 与 Hamming 核}

当 \texttt{feature\_types()} 包含任一 \texttt{"categorical"} 或 \texttt{"ordinal"} 标签时，\texttt{SurrogateModel} 自动切换到 BoTorch \texttt{MixedSingleTaskGP}\cite{ru2020cocabo}：
\begin{equation}
    k_{\text{mix}}(x, x') = k_{\text{Mat\'{e}rn}}\bigl(x^{\text{cont}}, x'^{\text{cont}}\bigr) \cdot k_{\text{cat}}\bigl(x^{\text{cat}}, x'^{\text{cat}}\bigr)
\end{equation}
其中 $k_{\text{cat}}$ 为 Hamming 核的连续松弛 \texttt{CategoricalKernel}。由于当前 CO$_2$RR 未声明 \texttt{categorical} 维度，此路径处于待激活状态；当老 BoTorch 版本不支持 \texttt{MixedSingleTaskGP} 时，\texttt{SurrogateModel.fit} 优雅降级为 \texttt{SingleTaskGP} 并打印告警（仅影响未来可能的类别化 Task）。

\subsection{动态候选池生成}

\texttt{TaskBase} 新增可选 Hook \texttt{generate\_candidates(n, seed)}：若 Task 实现之，orchestrator 在 Phase 3 入口\textbf{优先使用动态生成器替代静态 \texttt{candidates.csv}}。\texttt{CO2RRTask.generate\_candidates} 的实现：
\begin{itemize}
    \item 连续维：\texttt{torch.quasirandom.SobolEngine}（准随机序列，空间覆盖优于均匀 i.i.d.）
    \item 整数维：\texttt{numpy.random.Generator.integers}（在 $[l_k, h_k]$ 内均匀随机整数）
    \item 默认 $n = 1000$，由 \texttt{-{}-generate-candidates-n} 控制
\end{itemize}

回退机制三重冗余：(1) Task 未实现生成器 $\to$ 回退 CSV；(2) 生成抛异常 $\to$ 回退 CSV；(3) 用户显式 \texttt{-{}-prefer-file-candidates} $\to$ 强制 CSV（复现旧实验、无入侵修改）。

\subsection{Thompson 离散策略}

当候选池规模达到数千时，对每个候选重复计算采集函数代价昂贵，且 UCB 值差异可能集中在局部导致 Top-N 推荐同质化。新增 \texttt{-{}-discrete-strategy thompson} 采用 Thompson Sampling\cite{kandasamy2018parallel}：从 GP 后验独立抽 $4q$ 个样本，每个样本上取候选池 argmax 作为 $q$ 之一，去重后截至 $q$ 个。
\begin{equation}
    x^{(j)} = \arg\max_{x \in D_{\text{cand}}} f^{(j)}(x), \quad f^{(j)} \sim p\bigl(f \mid \mathcal{D}_t\bigr), \quad j = 1, \ldots, 4q
\end{equation}
此举在 $O(|D_{\text{cand}}| \cdot q)$ GP forward 代价下实现\textbf{天然探索多样性}，无需额外 diversity regularizer（§8 的贪心子模多样性仅在 \texttt{acq} 策略下启用）。失败回退：抽样异常则退化为按后验均值排序 Top-N。

\subsection{架构合规性与体系无关性校验}

所有改动严格遵循"算法层零体系常量"原则：
\begin{itemize}
    \item \texttt{kabo/utils.py} 新增 \texttt{round\_integer\_dims\_to\_grid} / \texttt{integer\_indices\_from\_types} / \texttt{categorical\_indices\_from\_types}，\textbf{仅接受 tensor 与索引 list}
    \item \texttt{kabo/surrogate.py::SurrogateModel.fit} 按 \texttt{feature\_types} 字典分派，不硬编码 "hbond" 等词
    \item \texttt{kabo/acquisition.py::optimize\_continuous} 新增 \texttt{integer\_indices} 纯索引参数
    \item \texttt{kabo/engine.py::KABOEngine} 新增 \texttt{discrete\_strategy} 字符串开关
\end{itemize}

对 \texttt{kabo/engine.py}、\texttt{kabo/acquisition.py}、\texttt{kabo/surrogate.py} 三个算法层文件执行 \texttt{grep "hbond\textbar{}CO2\textbar{}pI"} 扫描均无匹配，合规性通过静态校验。

\subsection{端到端验证}

在五种配置组合下端到端回归全部通过：
\begin{enumerate}
    \item CO$_2$RR 默认（Sobol 生成 + \texttt{acq}）
    \item CO$_2$RR Sobol 生成 + \texttt{thompson}
    \item CO$_2$RR \texttt{-{}-prefer-file-candidates}（legacy CSV + \texttt{acq}）
    \item TestTask（纯连续，CSV 回退路径，向后兼容性基线）
    \item CO$_2$RR KABO mode + 2 iterations（端到端含专家先验与偏好学习）
\end{enumerate}

代表性证据：
\begin{itemize}
    \item CO$_2$RR 所有推荐配方中 6 个氢键维均为整数（如 \texttt{A\_\brk hbond\_\brk acceptors = 3.0000}、\texttt{B\_\brk hbond\_\brk donors = 2.0000}），连续维仍为 Sobol 浮点（如 \texttt{A\_\brk pI = 3.4407}）
    \item TestTask 未声明整数维 $\to$ UCB 值 2.0227 与实施前逐字节一致，向后兼容零回归
    \item Thompson 策略日志 "drew 20 samples, returned 5 unique picks"，多样性可审计
\end{itemize}


\section{知识增强贝叶斯优化 (KABO)}

标准 HITL-BO 中，专家反馈仅体现为单次候选选择，其蕴含的偏好模式和领域知识未被结构性沉淀。KABO 通过三条知识通道注入结构化先验。

\subsection{KABO 复合采集函数}

KABO 核心是一个四项加权组合采集函数。为防止各项因量纲与数值范围差异导致的信号压制，在组合前对每一项做\textbf{在线 z-score 标准化}：

\begin{equation}
    \alpha_{\text{KABO}}(x) = z\bigl(\alpha_{\text{base}}(x)\bigr)
    + \lambda_p \cdot z\bigl(\text{Pref}_{\text{MC}}(x)\bigr)
    + \lambda_k \cdot z\bigl(\text{Prior}(x)\bigr)
    + \lambda_v \cdot z\bigl(\text{VOI}(x)\bigr)
    \label{eq:kabo}
\end{equation}

z-score 定义为：

\begin{equation}
    z(t) = \frac{t - \bar{t}}{\hat{\sigma}_t + \epsilon}, \quad \hat{\sigma}_t = \sqrt{\tfrac{1}{n} \sum_{i=1}^n (t_i - \bar{t})^2}, \quad \epsilon = 10^{-8}
\end{equation}

批次为单元素或标准差为零时 $z(\cdot)$ 返回全零以避免数值退化。各项含义：

\begin{itemize}
    \item $\alpha_{\text{base}}$：基础采集（UCB 或 qNEI），代理驱动的探索-开发平衡
    \item $\text{Pref}_{\text{MC}}$：偏好后验的 MC 积分得分（\S\ref{sec:pref}）
    \item $\text{Prior}$：基于 JSON 配置的专家先验打分（\S\ref{sec:prior}）
    \item $\text{VOI}$：代理模型后验方差作为信息价值近似（\S\ref{sec:voi}）
\end{itemize}

权重 $\lambda_p, \lambda_k, \lambda_v \geq 0$ 均可 CLI 配置，默认 $\lambda_p = \lambda_k = 1.0$，$\lambda_v = 0.0$（关闭）。

\subsection{偏好学习与蒙特卡洛效用积分}
\label{sec:pref}

专家每轮选择执行某一候选而非其余，构成一组典型的成对比较数据。系统使用 BoTorch 的 \texttt{PairwiseGP}\cite{botorch2019} 以 Bradley-Terry 模型建模专家效用函数 $u(x)$：

\begin{equation}
    P(x_i \succ x_j) = \Phi\bigl(u(x_i) - u(x_j)\bigr), \quad u \sim \mathcal{GP}(\mathbf{0}, k_{\text{pref}})
\end{equation}

$\Phi$ 为标准正态 CDF，$k_{\text{pref}}$ 为偏好空间上的 ARD Mat\'{e}rn 核。

\textbf{关键改进（EI-UU）}：依 Astudillo \& Frazier\cite{astudillo2019} 的 Expected Improvement under Utility Uncertainty 框架，不使用偏好后验均值作为点估计，而是对效用后验做 MC 积分：

\begin{equation}
    \text{Pref}_{\text{MC}}(x) = \frac{1}{S} \sum_{s=1}^{S} u^{(s)}(x), \quad u^{(s)} \sim p(u | \mathcal{D}_{\text{pref}})
\end{equation}

默认 $S = 16$ 个后验采样，通过重参数化 (\texttt{rsample}) 实现。该方法保留效用不确定性信息，避免过早锁定偏差估计。

\textbf{偏好数据去重}：为防止候选重复导致比较图冗余膨胀，新点写入前做 $L_2$ 近邻查重（阈值 $\epsilon = 10^{-4}$），匹配到已有节点时复用索引。

\textbf{三态偏好支持}：交互界面支持 winner / loser / \textbf{tie} 三种回答。Tie 不写入偏好模型（\texttt{PairwiseGP} 不原生支持等价标签），仅作审计记录，\texttt{tie\_count} 写入实验元数据。

\subsection{专家先验编码}
\label{sec:prior}

领域专家对某些描述符具有先验判断（例如"最优电位约 $-1.0$ 至 $-0.6$ V"）。系统通过 JSON 配置将这些知识编码为参数化分布（高斯或均匀），采集评估时产生先验打分：

\begin{equation}
    \text{Prior}(x) = \sum_{k \in \mathcal{K}} S_k(x_k)
\end{equation}

$\mathcal{K}$ 为配置了先验的特征子集。高斯先验 $S_k(x_k) = -\frac{(x_k - \mu_k)^2}{2\sigma_k^2}$（对数密度核），均匀先验在区间外施加常数惩罚。

\textbf{先验校验流程 (Prior Predictive Check)}：\texttt{scripts/validate\_prior.py --task <name>} 对 JSON 配置执行：(1) 参数合法性校验（$\sigma > 0$、$\min < \max$）；(2) Monte Carlo 采样检查设计空间覆盖率；(3) 输出结构化 JSON 报告。未通过校验的先验不应进入正式实验。每条记录含 \texttt{confidence}、\texttt{evidence} 审计字段保证可追溯。工具通过 \texttt{--task} 参数从 \texttt{TASK\_REGISTRY} 取得对应设计域边界，天然支持多体系。

\subsection{近似信息价值 (Approximate VOI)}
\label{sec:voi}

知识梯度 (Knowledge Gradient, KG)\cite{wu2016kg} 的核心是偏好"信息价值高"的候选。真正的 KG 需要 one-step lookahead 嵌套优化，计算开销极大。作为轻量近似，使用代理模型后验方差作为信息价值代理：

\begin{equation}
    \text{VOI}(x) = \text{Var}[f(x) | \mathcal{D}_t] = \sigma_{t}^2(x)
\end{equation}

高方差区域的评估最大化后验信息增益，这是 KG 的零阶近似。通过 $\lambda_v$ 控制权重，默认关闭。


\section{偏好探索子循环 (Preference Exploration)}

标准 HITL-BO 中偏好信息仅从实验后的候选选择中被动获取。受 Lin et al.\cite{lin2022pebo} PEBO 框架启发，系统引入\textbf{偏好探索 (Preference Exploration, PE)} 子循环：在每轮实物实验之前，以零实验代价主动查询专家偏好。

采用\textbf{不确定性策略}：从离散候选池中选择偏好模型最不确定、且预测偏好差异最小的 pair——信息量最大。对每对 $(x_i, x_j)$ 计算：

\begin{equation}
    \text{Info}(x_i, x_j) = \text{Var}[u(x_i)] + \text{Var}[u(x_j)] - \bigl|\mathbb{E}[u(x_i)] - \mathbb{E}[u(x_j)]\bigr|
\end{equation}

得分最高的 $B$ 对（$B = $ \texttt{pe\_budget}，默认 0 关闭）依次呈现给专家，每对支持 "A 优 / B 优 / Tie" 三种回答。非 tie 直接写入偏好模型，PE 结束后触发偏好 GP 重拟合，再进入正式采集优化。PE 在冷启动时（无偏好模型）退化为随机对选择。


\section{多样性感知推荐菜单}

受 Astudillo \& Frazier\cite{astudillo2019} "向决策者呈现菜单而非单点"理念驱动，Top-N 推荐采用贪心子模选择进行多样性重排：

\begin{enumerate}
    \item 第 1 位：取纯最高采集值候选
    \item 后续位：取使 $\hat{\alpha}(x_i) + w_d \cdot \min_{j \in \mathcal{S}} \|x_i - x_j\|_2$ 最大者，其中 $\mathcal{S}$ 为已选集，$w_d$ 为多样性权重（\texttt{--diversity-weight}，默认 0.5）
\end{enumerate}

这避免了同质化推荐，$w_d = 0$ 退化为纯分数排序，$w_d$ 增大则强化候选间的空间分散。


\section{人类在环机制与 KABO 执行流}

连续边界寻优得到的连续理想配方点 $x^*_{\text{cont}}$ 与读取外部历史离散库得到的候选最佳点 $X^*_{\text{disc}}$ 之间，模型未必保证它们对应的物化配置绝对可以在实验室合规组装。此时人工干预（Manual Override）与偏好探索作为控制流汇聚核心介入。

整体流程分为\textbf{主循环骨架}（算法 \ref{alg:kabo-outer}）与\textbf{单轮内层迭代}（算法 \ref{alg:kabo-inner}）两部分，后者为前者 FOR 循环体中被调用的子过程 \textsc{InnerIteration}$(t, \mathcal{T}, \mathcal{D}_{t-1}, D_{\text{cand}}, \mathcal{M}_{\text{pref}})$。

\begin{algorithm}[H]
\caption{KABO 主循环骨架（适用于任意 \texttt{TaskBase} 子类）}
\label{alg:kabo-outer}
\begin{algorithmic}[1]
\REQUIRE Task 对象 $\mathcal{T}$，总迭代数 $T$，初始观测组 $\mathcal{D}_0$，候选数据库 $D_{\text{cand}}$
\REQUIRE KABO 参数 $\lambda_p, \lambda_k, \lambda_v$；PE 预算 $B$；先验配置 $\mathcal{P}$
\STATE \textbf{Initialize}: $d \leftarrow |\mathcal{T}.\text{feature\_columns}()|$；默认目标取 $\mathcal{T}.\text{default\_target}()$，用户可覆写采集策略
\STATE \textbf{Phase 1}: 用 RF 回归器从 $\mathcal{D}_0$ 对目标拟合，排序筛选最显著的 $K$ 项作 $X_{\text{selected}}$
\STATE 加载专家先验 $\mathcal{P}$ 并执行 Prior Predictive Check 校验
\STATE 初始化偏好模型 $\mathcal{M}_{\text{pref}} \leftarrow \varnothing$
\FOR{$t = 1$ \textbf{to} $T$}
    \STATE $(x_{\text{final}}, \mathbf{Y}_{\text{all}}, \mathcal{M}_{\text{pref}}) \leftarrow \textsc{InnerIteration}(t, \mathcal{T}, \mathcal{D}_{t-1}, D_{\text{cand}}, \mathcal{M}_{\text{pref}})$
    \STATE $\mathcal{D}_t \leftarrow \mathcal{D}_{t-1} \cup (x_{\text{final}}, \mathbf{Y}_{\text{all}})$
\ENDFOR
\end{algorithmic}
\end{algorithm}

\begin{algorithm}[H]
\caption{\textsc{InnerIteration}：单轮 KABO 迭代}
\label{alg:kabo-inner}
\begin{algorithmic}[1]
\REQUIRE 迭代轮数 $t$，Task $\mathcal{T}$，历史数据 $\mathcal{D}_{t-1}$，候选库 $D_{\text{cand}}$，偏好模型 $\mathcal{M}_{\text{pref}}$
\STATE \textit{// \textbf{Step 0: Preference Exploration（可选）}}
\IF{KABO 模式启用 \textbf{and} $B > 0$}
    \STATE 从 $D_{\text{cand}}$ 选出 $B$ 个最具信息量的候选对
    \FOR{$b = 1$ \textbf{to} $B$}
        \STATE 向专家展示 pair，收集回答 $\in \{A, B, \text{Tie}\}$
        \IF{非 Tie}
            \STATE 将 winner/loser 写入 $\mathcal{M}_{\text{pref}}$
        \ENDIF
    \ENDFOR
    \STATE 重拟合 $\mathcal{M}_{\text{pref}}$（PairwiseGP Laplace MLL）
\ENDIF
\STATE \textit{// \textbf{Step 1: 构建 KABO 采集函数}}
\STATE $\mathcal{M}_{GP} \leftarrow \text{Fit}(\mathcal{D}_{t-1}, k_{\text{ARD}})$ \textit{// 由 Engine.fit\_surrogate 完成}
\STATE $\alpha_{\text{base}} \leftarrow$ UCB$(\beta_t)$ 或 qNEI
\STATE $\alpha_{\text{KABO}} \leftarrow z(\alpha_{\text{base}}) + \lambda_p z(\text{Pref}_{\text{MC}}) + \lambda_k z(\text{Prior}) + \lambda_v z(\text{VOI})$
\STATE \textit{// \textbf{Step 2: 候选优化与多样性推荐}}
\STATE $x_{\text{cont}} \leftarrow \arg\max_{x \in [0,1]^K} \alpha_{\text{KABO}}(x)$
\STATE 评估离散候选 $\hat{Y}_{\text{disc}} \leftarrow \text{Evaluate}(D_{\text{cand}}, \alpha_{\text{KABO}})$
\STATE 按多样性重排推举 Top-N 菜单（贪心子模）
\STATE Pre-fill-before-choice：未被选中的 $d-K$ 维用专家值或中点补齐
\STATE \textit{// \textbf{Step 3: Human Review}}
\IF{专家选择 \texttt{tie}}
    \STATE $x_{\text{final}} \leftarrow$ Rank \#1；跳过偏好记录
\ELSIF{专家选择 \texttt{manual override}}
    \STATE $x_{\text{final}} \leftarrow x_{\text{manual}}$；$(x_{\text{manual}} \succ \text{all Top-N})$ 写入 $\mathcal{M}_{\text{pref}}$
\ELSE
    \STATE $x_{\text{final}} \leftarrow x_{\text{chosen}}$；写入偏好对
\ENDIF
\STATE \textit{// \textbf{Step 4: Laboratory Execution}}
\STATE $\mathbf{Y}_{\text{all}} \leftarrow \mathcal{T}.\text{prompt\_observation}(\cdot)$ \textit{// 或 simulate\_observation}
\STATE \textbf{return} $(x_{\text{final}}, \mathbf{Y}_{\text{all}}, \mathcal{M}_{\text{pref}})$
\end{algorithmic}
\end{algorithm}

值得注意的是两个伪代码中的所有步骤\textbf{均不涉及体系特定术语}：无论 $\mathcal{T}$ 是 \texttt{CO2RRTask}、\texttt{TestTask} 还是未来的 \texttt{OERTask}、\texttt{NRRTask}，控制流完全一致，体系差异仅在 $\mathcal{T}$ 的方法体内表达。


\section{案例研究：光催化 CO$_2$ 还原 (\texttt{CO2RRTask})}

\subsection{\texorpdfstring{19 描述符 $\times$ 7 产物方案}{19 描述符 x 7 产物方案}}

\texttt{CO2RRTask} 在 \texttt{kabo/task/co2rr.py} 中实现，约 150 行。其特征维度按物化分组：

\begin{itemize}
    \item \textbf{氨基酸 A}（4 维）：A\_pI、A\_distance、A\_hbond\_acceptors、A\_hbond\_donors
    \item \textbf{氨基酸 B}（4 维）：B\_pI、B\_distance、B\_hbond\_acceptors、B\_hbond\_donors
    \item \textbf{卟啉 MOF}（2 维）：MOF\_potential、M\_CO\_binding\_energy
    \item \textbf{光敏剂}（2 维）：PS\_absorption\_wavelength、PS\_potential
    \item \textbf{溶剂}（4 维）：Solvent\_dielectric、Solvent\_hbond\_acceptors / donors、CO$_2$\_solubility
    \item \textbf{反应条件}（3 维）：H$_2$O\_\brk concentration、Sacrificial\_\brk agent\_\brk potential、Sacrificial\_\brk agent\_\brk concentration
\end{itemize}

产物列包含 CO、HCOOH、CH$_4$、C$_2$H$_4$、CH$_3$OH、C$_2$H$_5$OH、H$_2$，其中 H$_2$ 被识别为竞争性 HER 副反应。

\subsection{HER 选择性复合目标}

\texttt{CO2RRTask.build\_training\_target} 支持 H$_2$ 惩罚项：当 \texttt{h2\_penalty\_weight} $> 0$ 时，构造 $y_{\text{composite}} = Y_{\text{target}} - w \cdot Y_{H_2}$ 作为代理训练目标，引导优化器偏离 HER 主导区。该机制是 CO$_2$RR 专有，OER 等其他体系无需暴露——验证了 Task 层封装体系特殊逻辑的合理性。

\subsection{产物模拟分布（Demo 模式）}

\texttt{CO2RRTask.simulate\_observation} 根据 CO$_2$RR 产物特性设定：目标产物以 GP 训练均值为中心、$0.3\sigma$ 为幅度采样；H$_2$ 在 $[2, 15]$ 均匀分布（HER 基线范围）；其他产物服从指数分布（scale = 2.0）模拟痕量副产物。该分布是 CO$_2$RR 领域的经验先验，\textbf{不被 Engine 所感知}。


\section{体系无关性验证：\texttt{TestTask}}

为独立验证 KABO 架构的任务解耦性，系统内置最小综合任务 \texttt{TestTask}（位于 \texttt{kabo/task/\brk test\_task.py}，约 60 行）：

\begin{itemize}
    \item 特征维：3（\texttt{x1}、\texttt{x2}、\texttt{x3}），单位边界 $[0, 1]^3$
    \item 产物列：1（\texttt{y}）
    \item 训练目标：直接取 \texttt{y}，无复合项
    \item 模拟分布：以 GP 均值为中心、$0.5\sigma$ 采样单目标
\end{itemize}

合成数据使用形如 $y = 10 - 5 \sum_k (x_k - \mu_k^*)^2 + \epsilon$ 的平滑单峰函数（$\mu^* = (0.3, 0.7, 0.5)$）。通过命令

\begin{verbatim}
python -m kabo --task test --data data/test_data.csv --candidates none \
    --skip-feature-selection --non-interactive --iterations 2 --seed 42
\end{verbatim}

可验证：(1) 相同 \texttt{KABOEngine} 实例在 3 维空间上正确拟合；(2) Phase 3 迭代推荐收敛至真实极值附近（GP 学到的 ARD 长度尺度约 $0.28 / 0.39 / 1.54$，正确反映 $x_3$ 对目标响应的相对弱相关）；(3) 全程无体系硬编码泄漏。


\section{可追溯性与参数审计}

为保证实验结果完全可复现，系统在每次运行结束将所有超参数与配置写入结构化元数据文件 \texttt{run\_metadata.json}：

\begin{itemize}
    \item \textbf{基础参数}：\texttt{seed}、\texttt{beta\_schedule}、\texttt{beta\_trace}、\texttt{acq\_strategy}
    \item \textbf{KABO 参数}：\texttt{kabo\_mode}、\texttt{lambda\_p}、\texttt{lambda\_k}、\texttt{lambda\_v}
    \item \textbf{交互参数}：\texttt{diversity\_weight}、\texttt{pe\_budget}、\texttt{tie\_count}
    \item \textbf{先验参数}：\texttt{expert\_prior\_file}（含版本追溯）
    \item \textbf{任务参数}：\texttt{task}（记录 \texttt{task\_name()}）、\texttt{selected\_features}（实际使用的 $K$ 个描述符）、\texttt{feature\_types}（完整的维度类型声明，见 §\ref{sec:mixed-space}）
    \item \textbf{混合空间参数}（§\ref{sec:mixed-space}）：\texttt{discrete\_strategy}（\texttt{acq} / \texttt{thompson}）、\texttt{generate\_candidates\_n}（动态候选池规模）、\texttt{prefer\_file\_candidates}（是否强制 CSV）
\end{itemize}

所有参数均在优化器初始化阶段执行严格防呆校验（有限性、非负性、类型检查），异常值在运行前即被拦截。


\section{讨论与展望}

\subsection{任务解耦的工程价值}

将 CO$_2$RR 专用代码重构为任务解耦架构的核心收益体现在三方面：

\begin{enumerate}
    \item \textbf{可扩展性}：新增体系的增量代码量 $\sim 50$ 行，零算法改动；多组实验室可并行各自体系而不互相干扰。
    \item \textbf{可验证性}：\texttt{TestTask} 作为合成基线使 Engine 数学逻辑可以在合成数据上独立单元测试，CO$_2$RR 回归测试可以冻结基线比对。
    \item \textbf{可维护性}：体系常量、交互文案、模拟分布隔离在各自 Task 文件内，修改某个体系不会波及其他体系的测试基线。
\end{enumerate}

\subsection{潜在扩展方向}

\begin{itemize}
    \item \textbf{多目标 Pareto 前沿}：目前 Task 返回标量训练目标，未来可扩展 \texttt{build\_training\_targets()} 返回向量以支持 multi-objective qEHVI\cite{botorch2019}
    \item \textbf{异构特征类型（部分已实现，见 §\ref{sec:mixed-space}）}：\texttt{TaskBase.feature\_types()} 契约已支持连续、整数、类别、序数四类声明；\texttt{SurrogateModel} 对类别维自动切换 \texttt{MixedSingleTaskGP}\cite{ru2020cocabo}，对整数维应用 Round-trick\cite{garrido2020round}。下一步可将氨基酸 / 溶剂从连续描述符代理升级为显式类别，并配合 ESM2 / AAindex 化学嵌入作为类别先验核\cite{catbox2025}
    \item \textbf{更深的先验注入}：Mikkola et al.\cite{mikkola2021prior} 的先验启发方法可与 JSON 先验编码融合，支持领域专家以自然语言描述先验
    \item \textbf{高维 TuRBO 信任域}：当 OER 等任务的描述符维度 $d > 30$ 时引入 TuRBO\cite{turbo2019} 信任域切换，避免全局 GP 在高维下的病态拟合
\end{itemize}


\section{结语}

本报告构建了从特征解耦裁剪、物理感知映射的高斯代理，到知识增强贝叶斯优化闭环的完整系统，并通过任务解耦架构使同一套算法内核支持任意催化体系。相较于原始 HITL-BO 框架，KABO 在四个方面实现了质的提升：(1) 通过 PairwiseGP 偏好学习与 MC 积分将专家交互选择转化为持久化的效用模型；(2) 通过 JSON 参数化先验编码与 Prior Predictive Check 使领域知识的注入可审计、可校验；(3) 通过偏好探索子循环以零实验代价主动增强偏好模型质量；(4) 通过近似 VOI 项引导探索高信息价值区域。四项复合采集函数经在线 z-score 标定后协同工作，在 CO$_2$RR 等不确定大搜索面下兼顾开发、探索、偏好与知识价值四重目标。任务解耦层使这套机制得以从 CO$_2$RR 顺利推广至 OER、NRR、ORR 等同类仿酶催化体系，成为多成分催化实验与自动化合成间的可迁移驱动支点。

\newpage
\begin{thebibliography}{99}

\bibitem{gpucb2026}
Ding et al., \textit{Efficient and Principled Scientific Discovery through Bayesian Optimization: A Tutorial}. arXiv:2604.01328v3.

\bibitem{cbo2017}
Letham et al., \textit{Constrained Bayesian Optimization with Noisy Experiments}. arXiv:1706.07094v2.

\bibitem{botorch2019}
Balandat et al., \textit{BoTorch: A Framework for Efficient Monte-Carlo Bayesian Optimization}. arXiv:1910.06403v3.

\bibitem{turbo2019}
Wang et al., \textit{Scalable Global Optimization via Local Bayesian Optimization}. arXiv:1910.01739v4.

\bibitem{bo2012}
Snoek et al., \textit{Practical Bayesian Optimization of Machine Learning Algorithms}. arXiv:1206.2944v2.

\bibitem{catbox2025}
Li et al., \textit{CatBOX: A Categorical-Continuous Bayesian Optimization with Spectral Mixture Kernels for Accelerated Catalysis Experiments}. arXiv:2505.17393v2.

\bibitem{jacs2024}
Author et al., \textit{Multivariate Bayesian Optimization of CoO Nanoparticles for CO2 Hydrogenation Catalysis}. Journal of the American Chemical Society (2024), DOI: 10.1021/jacs.4c03789.

\bibitem{nature2025}
Author et al., \textit{Identifying a highly efficient molecular photocatalytic CO2 reduction system via descriptor-based high-throughput screening}. Nature Catalysis (2025), DOI: 10.1038/s41929-025-01291-z.

\bibitem{astudillo2019}
Astudillo \& Frazier, \textit{Multi-Attribute Bayesian Optimization under Utility Uncertainty}. arXiv:1911.05934.

\bibitem{wu2016kg}
Wu \& Frazier, \textit{The Parallel Knowledge Gradient Method for Batch Bayesian Optimization}. arXiv:1606.04414.

\bibitem{li2020prior}
Li et al., \textit{Incorporating Expert Prior in Bayesian Optimisation via Space Warping}. arXiv:2002.11256.

\bibitem{lin2022pebo}
Lin et al., \textit{Preference Exploration for Efficient Bayesian Optimization with Multiple Outcomes}. arXiv:2203.11382.

\bibitem{mikkola2021prior}
Mikkola et al., \textit{Prior Knowledge Elicitation: The Past, Present, and Future}. arXiv:2112.01380.

\bibitem{garrido2020round}
Garrido-Merch\'{a}n E C \& Hern\'{a}ndez-Lobato D, \textit{Dealing with Categorical and Integer-valued Variables in Bayesian Optimization with Gaussian Processes}. Neurocomputing 380 (2020) 20--35. arXiv:1805.03463.

\bibitem{ru2020cocabo}
Ru B, Alvi A, Nguyen V, Osborne M A, Roberts S J, \textit{Bayesian Optimisation over Multiple Continuous and Categorical Inputs}. ICML 2020. arXiv:1906.08878.

\bibitem{kandasamy2018parallel}
Kandasamy K, Krishnamurthy A, Schneider J, P\'{o}czos B, \textit{Parallelised Bayesian Optimisation via Thompson Sampling}. AISTATS 2018. arXiv:1705.09236.

\end{thebibliography}

\end{document}
